\begin{figure}[htbp]
    \centering
    \begin{tikzpicture}[
        block/.style={draw, rectangle, minimum height=1.2cm, minimum width=2.6cm, rounded corners=3pt, line width=0.8pt, align=center, fill=white, font=\small\sffamily\bfseries},
        nobox/.style={rectangle, minimum height=1.2cm, minimum width=2.6cm, align=center, font=\small\sffamily\bfseries},
        arrow/.style={-{Stealth[scale=1.2]}, line width=0.8pt}
    ]
    
    % Main Horizontal Chain
    \node (raster) [nobox, minimum width =1.5cm] {CSR};
    \node (vcn) [block, right=0.8cm of raster] {VCN/VNLL\\Model};
    \node (ic) [block, minimum width =1.5cm, right=1.5cm of vcn] {IC};
    \node (ac) [block, minimum width =1.5cm, right=0.8cm of ic] {AC};
    \node (sc) [block, minimum width =1.5cm, right=0.8cm of ac] {SC};
    \node (readout) [nobox, right=0.8cm of sc] {Distance\\Readout};
    
    % Detour Node (Below and slightly right of VCN)
    \node (dnll) [block, below=1.0cm of vcn, xshift=1.8cm] {3) DNLL};
    \node (cd) [block, minimum width =1.5cm, below=1.0cm of ac] {CD};
    
    % Connections
    \draw [arrow] (raster) -- (vcn);
    \draw [arrow] (vcn) -- (ic); % Direct feed
    
    % Detour Routing
    \draw [arrow] (vcn.south) |- (dnll.west);
    \draw [arrow, dashed] (dnll.east) -| ([xshift=-0.2cm]ic.south);
    \draw [arrow] (cd.west) -| ([xshift=+0.2cm]ic.south);
    
    
    % Remaining Chain
    \draw [arrow] (ic) -- (ac);
    \draw [arrow] (ac) -- (sc);
    \draw [arrow] (sc) -- (readout);
    
    \end{tikzpicture}
    \caption{Distance mapping pipeline featuring parallel brainstem inhibitory pathways.}
    \label{fig:distance_flowchart}
\end{figure}
